\par\noindent\textbf{(1)} 圆柱高斯面半径r∈(a,b)给2πr εE\_r=λ，故E\_r=λ/(2πεr)；Ampère环路2πr Hφ=I，故Hφ=I/(2πr)。导体表面切向电场为零，横向E与H均为TEM解，能量沿导线轴向流动。
\par\noindent\textbf{(2)} 两导体电势差V=∫\_a\textasciicircum{}b E\_rdr=λ ln(b/a)/(2πε)，故C′=λ/V=2πε/ln(b/a)。磁通或磁能法取U\_B′=∫\_a\textasciicircum{}b μH²/2·2πrdr=μI² ln(b/a)/(4π)=L′I²/2，因此L′=μ ln(b/a)/(2π)。
\par\noindent\textbf{(3)} 电能每长U\_E′=C′V²/2，磁能U\_B′=L′I²/2。Poynting功率P=∫(E×H)\_z2πrdr=VI，而传播本征解有U\_E′=U\_B′，从而V/I=√(L′/C′)=Z0=[ln(b/a)/(2π)]√(μ/ε)，与常用TEM传输线阻抗一致。
\par\noindent\textbf{(4)} 速度v=1/√(L′C′)=1/√(με)，几何对数完全抵消；半径比只改变阻抗而不改变均匀材料中的波速。b→a时C′→∞、L′→0但乘积L′C′=με保持固定；此形式极限不适合实际有有限壁厚与击穿场约束的电缆。
